\documentclass[10pt,letterpaper]{article}
\usepackage[utf8]{inputenc}
\usepackage[vietnamese]{babel}
\usepackage{amsmath,amssymb,amsfonts}
\usepackage{geometry}
\geometry{
    top=1.8cm,
    bottom=2cm,
    left=1.5cm,
    right=1.5cm,
    headheight=14pt,
    headsep=12pt
}
\usepackage{xcolor}
\usepackage{tcolorbox}
\usepackage{fancyhdr}
\usepackage{newtxtext,newtxmath}

% Định nghĩa màu chuẩn theo phong cách Stewart Calculus
\definecolor{stewartcyan}{RGB}{0, 136, 206}
\definecolor{stewartred}{RGB}{204, 34, 41}

% Khung định nghĩa theo quy tắc
\newtcolorbox{definitionbox}[1][]{
    colback=white,
    colframe=stewartred,
    boxrule=0.8pt,
    sharp corners,
    boxsep=0pt,
    left=10pt,
    right=10pt,
    top=7pt,
    bottom=7pt,
    before skip=7pt,
    after skip=7pt,
    #1
}

% Cấu hình Header và Footer
\pagestyle{fancy}
\fancyhf{}
\fancyhead[L]{\fontsize{9}{11}\selectfont \textbf{110}\qquad \textbf{\textsf{CHƯƠNG 2}}\quad \textsf{Giới hạn và Đạo hàm}}
\fancyhead[R]{}
\fancyfoot[C]{}
\renewcommand{\headrulewidth}{0pt}

\setlength{\parindent}{14pt}
\setlength{\parskip}{2pt}

\begin{document}

\noindent
\begin{minipage}[t]{0.26\textwidth}
    % Cột lề trái: để trống tương ứng với bản gốc trang 110
    \mbox{}
\end{minipage}%
\hfill
\begin{minipage}[t]{0.71\textwidth}
    \vspace{-4pt}
    \noindent Để ý rằng nếu chúng ta có thể tìm được một hằng số dương $C$ sao cho $|x + 3| < C$, thì
    \[
    |x + 3||x - 3| < C|x - 3|
    \]
    và ta có thể làm cho $C|x - 3| < \varepsilon$ bằng cách chọn $|x - 3| < \varepsilon/C$, do đó ta có thể chọn $\delta = \varepsilon/C$.

    Chúng ta có thể tìm một số $C$ như vậy nếu ta giới hạn $x$ nằm trong một khoảng nào đó có tâm tại $3$. Trên thực tế, vì chúng ta chỉ quan tâm đến các giá trị của $x$ ở gần $3$, nên hoàn toàn hợp lý khi giả định rằng khoảng cách từ $x$ đến $3$ không quá $1$, nghĩa là $|x - 3| < 1$. Khi đó $2 < x < 4$, nên $5 < x + 3 < 7$. Như vậy ta có $|x + 3| < 7$, và do đó $C = 7$ là một lựa chọn phù hợp cho hằng số.

    Nhưng bây giờ có hai điều kiện ràng buộc đối với $|x - 3|$, cụ thể là
    \[
    |x - 3| < 1 \qquad \text{và} \qquad |x - 3| < \frac{\varepsilon}{C} = \frac{\varepsilon}{7}
    \]
    Để đảm bảo rằng cả hai bất đẳng thức này đều được thỏa mãn, chúng ta chọn $\delta$ là số nhỏ hơn trong hai số $1$ và $\varepsilon/7$. Ký hiệu cho điều này là $\delta = \min\{1, \varepsilon/7\}$.

    \textbf{2.}\, \textit{Chứng minh rằng giá trị $\delta$ này thỏa mãn.}\quad Cho $\varepsilon > 0$, đặt $\delta = \min\{1, \varepsilon/7\}$. Nếu $0 < |x - 3| < \delta$, thì $|x - 3| < 1 \Rightarrow 2 < x < 4 \Rightarrow |x + 3| < 7$ (như ở phần 1). Ta cũng có $|x - 3| < \varepsilon/7$, do đó
    \[
    |x^2 - 9| = |x + 3||x - 3| < 7 \cdot \frac{\varepsilon}{7} = \varepsilon
    \]
    Điều này chứng tỏ rằng $\lim_{x \to 3} x^2 = 9$. \hfill {\color{stewartcyan}\rule{5pt}{5pt}}

    \vspace{0.8em}
    \noindent{\color{stewartcyan}\rule[0pt]{6.5pt}{6.5pt}}\hspace{0.5em}\textbf{\textsf{Giới hạn một phía}}\quad Các định nghĩa trực quan về giới hạn một phía đã được đưa ra trong Mục 2.2 có thể được phát biểu lại một cách chính xác như sau.

    \begin{definitionbox}
        \noindent{\color{stewartred}\setlength{\fboxrule}{0.8pt}\setlength{\fboxsep}{2pt}\fbox{\bfseries\sffamily\color{stewartred}3}}\hspace{0.5em}\textbf{\textsf{\color{stewartred}Định nghĩa chính xác về giới hạn bên trái}}
        \[
        \lim_{x \to a^-} f(x) = L
        \]
        \noindent nếu với mọi số $\varepsilon > 0$ đều tồn tại một số $\delta > 0$ sao cho
        \[
        \text{nếu} \qquad a - \delta < x < a \qquad \text{thì} \qquad |f(x) - L| < \varepsilon
        \]
    \end{definitionbox}

    \begin{definitionbox}
        \noindent{\color{stewartred}\setlength{\fboxrule}{0.8pt}\setlength{\fboxsep}{2pt}\fbox{\bfseries\sffamily\color{stewartred}4}}\hspace{0.5em}\textbf{\textsf{\color{stewartred}Định nghĩa chính xác về giới hạn bên phải}}
        \[
        \lim_{x \to a^+} f(x) = L
        \]
        \noindent nếu với mọi số $\varepsilon > 0$ đều tồn tại một số $\delta > 0$ sao cho
        \[
        \text{nếu} \qquad a < x < a + \delta \qquad \text{thì} \qquad |f(x) - L| < \varepsilon
        \]
    \end{definitionbox}

    \vspace{0.3em}
    Để ý rằng Định nghĩa 3 tương tự như Định nghĩa 2 ngoại trừ việc $x$ bị giới hạn nằm trong nửa \textit{bên trái} $(a - \delta, a)$ của khoảng $(a - \delta, a + \delta)$. Trong Định nghĩa 4, $x$ bị giới hạn nằm trong nửa \textit{bên phải} $(a, a + \delta)$ của khoảng $(a - \delta, a + \delta)$.
\end{minipage}

\end{document}